\ifmostrarGabarito
\begin{minted}{python}
def cifrar_texto(mensagem: str, deslocamento: int) -> str:
    resultado = []
    for c in mensagem:
        if 'A' <= c <= 'Z':
            novo_ord = (ord(c) - ord('A') + deslocamento) % 26 + ord('A')
            resultado.append(chr(novo_ord))
        else:
            resultado.append(c)
    return "".join(resultado)
\end{minted}

\textbf{Explicação:}
Iteramos pelos caracteres da mensagem. Se o caractere for uma letra maiúscula ('A' a 'Z'), calculamos a posição relativa no alfabeto (0 a 25), somamos o deslocamento com aritmética modular (\texttt{\% 26}) e reconvertemos para caractere. Caracteres não alfabéticos são mantidos inalterados.
\else
\linhasResposta{18}
\clearpage
\linhasResposta[FOLHA DE CONTINUACAO / RASCUNHO -- QUESTAO Q10]{30}
\clearpage
\fi
